\documentclass[11pt]{article}
\usepackage[margin=1in]{geometry}
\usepackage{booktabs}
\usepackage{amsmath}
\usepackage{hyperref}
\usepackage{xcolor}
\usepackage{longtable}
\usepackage{enumitem}
\usepackage{fancyvrb}
\usepackage{graphicx}

\hypersetup{colorlinks=true,linkcolor=blue,urlcolor=blue}

\title{Carryover Predictor Comparison in pmsimstats2025:\\
  A \texttt{tsd} Computation Bug and the Original's\\
  Open Power Question}
\author{pmsimstats2025 Project Notes}
\date{March 2026}

\begin{document}
\maketitle

\section{Context}

The \texttt{carryover\_focus/} simulation isolates the effect of
pharmacological carryover on statistical power for biomarker moderation
detection in the Hybrid N-of-1 aggregated trial design
\cite{hendrickson2020}. The simulation uses Hendrickson's binary drug
exposure predictor \texttt{Dbc} (bounded $[0, 1]$), which equals~1
when on drug and decays exponentially when off drug.

This document reports two findings:

\begin{enumerate}
  \item The pmsimstats2025 reimplementation computes the
    time-since-discontinuation variable (\texttt{tsd}) differently
    from Hendrickson's original \texttt{buildtrialdesign()} function.
    The 2025 code produces $t_{\text{sd}} = 0$ at every first off-drug
    measurement, while the original produces $t_{\text{sd}} \geq 1$
    (the interval width leading into that measurement). This
    discrepancy causes a step-function power cliff in the 2025
    simulation that does not exist in the original framework.

  \item Independent of this bug, the original Hendrickson simulation
    shows power declining more rapidly with increasing carryover
    half-life than pharmacological reasoning would predict. The cause
    of this behavior remains open and requires investigation after the
    \texttt{tsd} computation is corrected to match the original.
\end{enumerate}

\section{The Two Analysis Models}

The simulation compares two model specifications, toggled by the
\texttt{model\_carryover} flag. Both follow Hendrickson et~al.'s
\cite{hendrickson2020} binary predictor approach:

\medskip
\noindent\textbf{Model~A} (\texttt{model\_carryover = TRUE}):
\begin{verbatim}
response ~ Dbc * bm_centered + week
\end{verbatim}

\noindent\textbf{Model~B} (\texttt{model\_carryover = FALSE}):
\begin{verbatim}
response ~ treatment * bm_centered + week
\end{verbatim}

Both are fitted with \texttt{nlme::lme()} using a random intercept
per participant and a continuous-time AR(1) correlation structure
(\texttt{corCAR1(form = \~{} week | participant\_id)}), which models
the decay of within-subject correlation as a function of the time gap
between measurements. The difference between the two models lies
entirely in the drug exposure predictor used in the interaction term.

\subsection{Model Variable Definitions}

\begin{description}[style=nextline, leftmargin=2em]
  \item[\texttt{response}]
    The simulated clinical outcome for each participant at each
    measurement week. In the DGP, it is the sum of baseline status,
    biological response to drug (BR), expectancy response (ER), and
    time-in-trial trend (TR):
    $\text{response} = \text{baseline} + \text{BR} + \text{ER}
    + \text{TR}$.

  \item[\texttt{week}]
    The calendar week of the measurement within the trial schedule
    (e.g., 4, 8, 9, 10, 11, 12, 16, 20 for the Hybrid design).
    Included in the model as a fixed-effect covariate to absorb the
    secular time trend.

  \item[\texttt{bm\_centered}]
    The participant's biomarker value, standardized to zero mean and
    unit variance across the sample:
    $\texttt{bm\_centered}_i
    = (\text{biomarker}_i - \bar{\text{biomarker}})
    / \text{SD}(\text{biomarker})$.
    The drug-by-biomarker interaction term tests whether participants
    with higher biomarker values respond differently to treatment.

  \item[\texttt{participant\_id}]
    Unique identifier for each participant ($i = 1, \ldots, 70$).
    Defines the grouping structure for the random intercept and the
    within-subject correlation.

  \item[\texttt{treatment}]
    Binary indicator: 1 if the participant is on drug at that
    measurement week, 0 if off drug. Determined by the Hybrid design's
    path assignment (4 paths with staggered discontinuation and
    reinitiation schedules). In Model~B, this serves as the ``naive''
    drug exposure predictor: it assumes zero carryover.
\end{description}

\subsection{The \texttt{Dbc} Predictor}

The drug exposure predictor \texttt{Dbc} (``drug binary with
carryover'') follows Hendrickson et~al.\ \cite{hendrickson2020}. It
is bounded in $[0, 1]$:
\[
  \texttt{Dbc} =
  \begin{cases}
    1 & \text{if on drug} \\[4pt]
    0 & \text{if off drug and } t_{1/2} = 0
      \quad \text{(no carryover)} \\[4pt]
    \left(\tfrac{1}{2}\right)^{t_{\text{sd}} / t_{1/2}}
      & \text{if off drug and } t_{1/2} > 0
  \end{cases}
\]
Here $t_{\text{sd}}$ is the time since discontinuation in weeks and
$t_{1/2}$ is the carryover half-life (the time for the residual drug
effect to decay to half its value at discontinuation). The critical
question is how $t_{\text{sd}}$ is computed; the original and the 2025
reimplementation differ on this point (Section~\ref{sec:bug}).

\section{Hendrickson's Original \texttt{tsd} Computation}
\label{sec:original}

\subsection{The \texttt{buildtrialdesign()} Algorithm}

In the original pmsimstats package, \texttt{tsd} is computed by
\texttt{buildtrialdesign()} (\texttt{buildtrialdesign.R}, lines
86--106). The algorithm accumulates off-drug interval widths:

\begin{verbatim}
t_wk <- c(4, 4, 1, 1, 1, 1, 4, 4)   # interval widths
tsd  <- t_wk * (od != 1)              # seed: width if off-drug
for (iTP in 2:nP) {
  if (od_intervals[iTP] > 0)          # on drug
    tod[iTP] <- tod[iTP] + tod[iTP-1]
  else                                 # off drug
    tsd[iTP] <- tsd[iTP] + tsd[iTP-1]
}
tsd <- everondrug * tsd
\end{verbatim}

The implicit assumption is that drug discontinuation occurs at the
\emph{start} of the interval, so by the measurement at the end of the
interval the participant has been off drug for the full interval
duration. The first off-drug measurement therefore has
$t_{\text{sd}} = t_{\text{wk}}$ (the interval width), never zero.

\subsection{Original \texttt{tsd} Values for All Four Paths}

The Hybrid design uses measurement weeks $(4, 8, 9, 10, 11, 12, 16,
20)$ with interval widths $t_{\text{wk}} = (4, 4, 1, 1, 1, 1, 4, 4)$.
The four \texttt{ondrug} vectors from the publication vignette:

\begin{verbatim}
Path A: c(1, 1, 1, 1, 0, 0, 1, 0)
Path B: c(1, 1, 1, 1, 0, 0, 0, 1)
Path C: c(1, 1, 1, 0, 0, 0, 1, 0)
Path D: c(1, 1, 1, 0, 0, 0, 0, 1)
\end{verbatim}

\noindent These correspond to pmsimstats2025 Paths 1--4 respectively.
Tracing the accumulation algorithm for each path yields:

\begin{table}[h]
\centering
\caption{Original \texttt{tsd} values (in weeks) for all four
  Hybrid design paths, computed by \texttt{buildtrialdesign()}.
  On-drug rows show \texttt{tsd}~$= 0$; off-drug rows show
  cumulative interval widths. The minimum off-drug \texttt{tsd}
  is~1 (clustered phase) or~4 (spaced phase).}
\label{tab:orig_tsd}
\begin{tabular}{@{}ccccccccc@{}}
\toprule
 & \multicolumn{8}{c}{Week} \\
\cmidrule(l){2-9}
Path & 4 & 8 & 9 & 10 & 11 & 12 & 16 & 20 \\
\midrule
A (ondrug) & 1 & 1 & 1 & 1 & 0 & 0 & 1 & 0 \\
A (tsd)    & 0 & 0 & 0 & 0 & 1 & 2 & 0 & 4 \\
\midrule
B (ondrug) & 1 & 1 & 1 & 1 & 0 & 0 & 0 & 1 \\
B (tsd)    & 0 & 0 & 0 & 0 & 1 & 2 & 6 & 0 \\
\midrule
C (ondrug) & 1 & 1 & 1 & 0 & 0 & 0 & 1 & 0 \\
C (tsd)    & 0 & 0 & 0 & 1 & 2 & 3 & 0 & 4 \\
\midrule
D (ondrug) & 1 & 1 & 1 & 0 & 0 & 0 & 0 & 1 \\
D (tsd)    & 0 & 0 & 0 & 1 & 2 & 3 & 7 & 0 \\
\bottomrule
\end{tabular}
\end{table}

\subsection{Original \texttt{Dbc} in the Analysis Model}

In \texttt{lme\_analysis.R} (lines 107--109), the original computes
\texttt{Dbc} as:

\begin{verbatim}
data.m2[Db==TRUE,  Dbc := 1]
data.m2[Db==FALSE, Dbc := (1/2)^(sf * tsd / t_half)]
\end{verbatim}

\noindent where \texttt{sf} is
\texttt{op\$carryover\_scalefactor}, which defaults to~1 (not~7 as
previously stated in an earlier version of this document). Note that
the data-generating process in \texttt{generateData()} uses a separate
\texttt{scalefactor} parameter that defaults to~2 (see
Section~\ref{sec:open}).

\subsection{Original \texttt{Dbc} at Boundary Measurements}

With the original \texttt{tsd} values, the first off-drug measurement
in the clustered phase has $t_{\text{sd}} = 1$~week:

\begin{table}[h]
\centering
\caption{Original \texttt{Dbc} at off-drug measurements for a path~C
  participant, using the original interval-based \texttt{tsd}.
  Boundary measurements carry negligible carryover for short
  half-lives.}
\label{tab:orig_dbc}
\begin{tabular}{@{}ccrr@{}}
\toprule
Week & Orig \texttt{tsd}
  & Orig \texttt{Dbc} ($t_{1/2} = 0.1$)
  & Orig \texttt{Dbc} ($t_{1/2} = 0.2$) \\
\midrule
10 & 1 & $(1/2)^{10} \approx 0.001$ & $(1/2)^{5} = 0.031$ \\
11 & 2 & $(1/2)^{20} \approx 0$ & $(1/2)^{10} \approx 0.001$ \\
12 & 3 & $(1/2)^{30} \approx 0$ & $(1/2)^{15} \approx 0$ \\
20 & 4 & $(1/2)^{40} \approx 0$ & $(1/2)^{20} \approx 0$ \\
\bottomrule
\end{tabular}
\end{table}

\noindent For $t_{1/2} = 0.1$~weeks, the boundary \texttt{Dbc} is
0.001---essentially zero carryover. Even at $t_{1/2} = 0.2$~weeks,
the boundary value is only 0.031. These values produce a smooth,
gradual decline in the on/off contrast as the half-life increases,
not a step function.

\section{The 2025 Reimplementation Discrepancy}
\label{sec:bug}

\subsection{The \texttt{carryover\_power\_simulation.R} Algorithm}

The 2025 code computes \texttt{tsd} as the difference between the
current calendar week and the week at which a treatment transition
(on $\to$ off) is detected:

\begin{verbatim}
just_discontinued = treatment == 0 &
  lag(treatment, default = 0) == 1
discontinuation_week = if_else(
  just_discontinued, week, NA_real_
)
discontinuation_week = zoo::na.locf(
  discontinuation_week, na.rm = FALSE
)
tsd = if_else(
  treatment == 0 & !is.na(discontinuation_week),
  pmax(week - discontinuation_week, local_tsd_min),
  0
)
\end{verbatim}

The implicit assumption is that discontinuation occurs \emph{at} the
measurement week. At the first off-drug measurement,
\texttt{week == discontinuation\_week}, so
$t_{\text{sd}} = 0$.

\subsection{Side-by-Side \texttt{tsd} Comparison}

Table~\ref{tab:tsd_compare} compares the original and 2025
\texttt{tsd} values for all four paths. Only off-drug rows are shown.

\begin{table}[h]
\centering
\caption{Original vs.\ 2025 \texttt{tsd} at off-drug measurements.
  The 2025 value is systematically one interval width shorter at the
  first off-drug measurement after each discontinuation. Boldface
  marks mismatches.}
\label{tab:tsd_compare}
\begin{tabular}{@{}ccccc@{}}
\toprule
Path & Week & Orig \texttt{tsd} & 2025 \texttt{tsd} & Match? \\
\midrule
A & 11 & 1 & \textbf{0} & No \\
A & 12 & 2 & \textbf{1} & No \\
A & 20 & 4 & \textbf{0} & No \\
\midrule
B & 11 & 1 & \textbf{0} & No \\
B & 12 & 2 & \textbf{1} & No \\
B & 16 & 6 & \textbf{5} & No \\
\midrule
C & 10 & 1 & \textbf{0} & No \\
C & 11 & 2 & \textbf{1} & No \\
C & 12 & 3 & \textbf{2} & No \\
C & 20 & 4 & \textbf{0} & No \\
\midrule
D & 10 & 1 & \textbf{0} & No \\
D & 11 & 2 & \textbf{1} & No \\
D & 12 & 3 & \textbf{2} & No \\
D & 16 & 7 & \textbf{6} & No \\
\bottomrule
\end{tabular}
\end{table}

\noindent Every off-drug measurement is affected. The offset equals
the interval width leading into the first off-drug timepoint: 1~week
in the clustered phase (weeks 10--12) and 4~weeks in the spaced phase
(weeks 16, 20). Re-discontinuation points (e.g., week~20 for paths~A
and~C) show the largest discrepancy: original \texttt{tsd}~$= 4$
versus 2025 \texttt{tsd}~$= 0$.

\subsection{Consequence for \texttt{Dbc}}

Table~\ref{tab:dbc_compare} shows the resulting \texttt{Dbc} values
for a path~C participant at $t_{1/2} = 0.1$~weeks.

\begin{table}[h]
\centering
\caption{\texttt{Dbc} under original vs.\ 2025 \texttt{tsd} for
  path~C, $t_{1/2} = 0.1$~weeks. The 2025 code produces full
  carryover (\texttt{Dbc}~$= 1$) at boundary measurements where the
  original produces negligible carryover.}
\label{tab:dbc_compare}
\begin{tabular}{@{}cccccc@{}}
\toprule
Week & Orig tsd & Orig Dbc & 2025 tsd & 2025 Dbc & Discrepancy \\
\midrule
10 & 1 & 0.001  & 0 & 1.000 & $\times$1000 \\
11 & 2 & $\approx$0 & 1 & 0.001 & --- \\
12 & 3 & $\approx$0 & 2 & $\approx$0 & --- \\
20 & 4 & $\approx$0 & 0 & 1.000 & $\infty$ \\
\bottomrule
\end{tabular}
\end{table}

At $t_{1/2} = 0.2$~weeks the pattern is the same:

\begin{table}[h]
\centering
\caption{\texttt{Dbc} under original vs.\ 2025 \texttt{tsd} for
  path~C, $t_{1/2} = 0.2$~weeks.}
\label{tab:dbc_compare_02}
\begin{tabular}{@{}cccccc@{}}
\toprule
Week & Orig tsd & Orig Dbc & 2025 tsd & 2025 Dbc & Discrepancy \\
\midrule
10 & 1 & 0.031  & 0 & 1.000 & $\times$32 \\
11 & 2 & 0.001  & 1 & 0.031 & $\times$31 \\
12 & 3 & $\approx$0 & 2 & 0.001 & --- \\
20 & 4 & $\approx$0 & 0 & 1.000 & $\infty$ \\
\bottomrule
\end{tabular}
\end{table}

\section{The Step-Function Power Cliff (a 2025 Bug)}
\label{sec:cliff}

\subsection{Mathematical Origin}

The step-function power cliff observed in the 2025 simulation arises
because $t_{\text{sd}} = 0$ at boundary measurements. The exponential
decay model evaluates as:
\[
  f(t_{1/2}) \big|_{t_{\text{sd}}=0} =
  \begin{cases}
    0 & \text{if } t_{1/2} = 0 \quad \text{(special case in code)} \\[4pt]
    \left(\dfrac{1}{2}\right)^{0 / t_{1/2}}
      = \left(\dfrac{1}{2}\right)^{0} = 1
      & \text{if } t_{1/2} > 0
  \end{cases}
\]

This is a step function: for \textbf{any} positive half-life, no
matter how small, the carryover effect at the boundary equals~1.
The parameter $t_{1/2}$ controls only the rate of subsequent decay,
not the initial value. This behavior is specific to the 2025
\texttt{tsd} computation and does not occur in the original, where
$t_{\text{sd}} \geq 1$ at all boundary measurements.

\subsection{Consequences for the 2025 Simulation}

In the 2025 simulation with \texttt{tsd\_min}~$= 0$ (uncorrected),
the naive model (\texttt{treatment}) shows a step-function collapse:

\begin{table}[h]
\centering
\caption{2025 simulation results with uncorrected \texttt{tsd}
  (\texttt{tsd\_min}~$= 0$, $N = 70$, 200 iterations). The naive
  model collapses to $\sim$0.06 power at any nonzero half-life.
  This behavior is caused by the \texttt{tsd} bug, not by a property
  of the carryover decay model.}
\label{tab:uncorrected}
\begin{tabular}{@{}cccccc@{}}
\toprule
$t_{1/2}$ (wks) & \multicolumn{2}{c}{Power} &
  \multicolumn{2}{c}{Mean $\hat{\beta}$} & Difference \\
\cmidrule(lr){2-3} \cmidrule(lr){4-5}
  & Dbc & Treatment & Dbc & Treatment & (power) \\
\midrule
0.00 & 0.73 & ---  & 0.675 & ---   & ---    \\
0.05 & 0.63 & 0.07 & 0.647 & 0.144 & +0.56  \\
0.10 & 0.69 & 0.06 & 0.694 & 0.133 & +0.64  \\
0.15 & 0.66 & 0.06 & 0.667 & 0.139 & +0.60  \\
0.20 & 0.57 & 0.07 & 0.631 & 0.124 & +0.50  \\
\bottomrule
\end{tabular}
\end{table}

The mechanism: at any nonzero $t_{1/2}$, weeks~10 and~20 (and their
equivalents across all four paths) carry $\texttt{Dbc} = 1$---full
drug effect---because the 2025 code assigns $t_{\text{sd}} = 0$. These
boundary measurements constitute half of all off-drug timepoints.
The naive model assigns \texttt{treatment}~$= 0$ to these timepoints
while they still carry full biomarker-moderated drug response,
collapsing the on/off contrast. The mean interaction estimate drops
from $\hat{\beta} \approx 0.68$ to $\hat{\beta} \approx 0.13$---an
80\% attenuation.

\subsection{The \texttt{tsd\_min} Workaround}

As part of the 2025 development, a minimum \texttt{tsd} floor was
introduced to address the step-function cliff:
\[
  t_{\text{sd,eff}} = \max(t_{\text{sd}}, \; t_{\text{sd,min}})
\]

With $t_{\text{sd,min}} = 0.5$~weeks, the boundary carryover becomes
a smooth function of $t_{1/2}$:

\begin{table}[h]
\centering
\caption{Boundary carryover under three regimes: the buggy 2025
  code (\texttt{tsd}~$= 0$), the \texttt{tsd\_min}~$= 0.5$
  workaround, and the original interval-based \texttt{tsd}~$= 1$
  (clustered phase).}
\label{tab:boundary_three}
\begin{tabular}{@{}cccc@{}}
\toprule
$t_{1/2}$ (wks)
  & 2025 bug ($t_{\text{sd}} = 0$)
  & Workaround ($t_{\text{sd}} = 0.5$)
  & Original ($t_{\text{sd}} = 1$) \\
\midrule
0.00 & 0 & 0 & 0 \\
0.05 & 1.000 & 0.001 & $\approx$0 \\
0.10 & 1.000 & 0.031 & 0.001 \\
0.15 & 1.000 & 0.101 & 0.010 \\
0.20 & 1.000 & 0.177 & 0.031 \\
\bottomrule
\end{tabular}
\end{table}

The workaround eliminates the step function and produces gradual
power decline:

\begin{table}[h]
\centering
\caption{2025 simulation results with \texttt{tsd\_min}~$= 0.5$
  ($N = 70$, 200 iterations). The step function is eliminated, but
  these results do not match the original's \texttt{tsd} computation,
  where the effective minimum is 1~week (clustered) or 4~weeks
  (spaced).}
\label{tab:corrected}
\begin{tabular}{@{}cccccc@{}}
\toprule
$t_{1/2}$ (wks) & \multicolumn{2}{c}{Power} &
  \multicolumn{2}{c}{Mean $\hat{\beta}$} & Difference \\
\cmidrule(lr){2-3} \cmidrule(lr){4-5}
  & Dbc & Treatment & Dbc & Treatment & (power) \\
\midrule
0.00 & 0.69 & ---  & 0.646 & ---   & ---    \\
0.05 & 0.68 & 0.72 & 0.663 & 0.669 & $-$0.04  \\
0.10 & 0.68 & 0.73 & 0.653 & 0.656 & $-$0.05  \\
0.15 & 0.66 & 0.58 & 0.653 & 0.579 & +0.08  \\
0.20 & 0.53 & 0.54 & 0.625 & 0.544 & $-$0.01  \\
\bottomrule
\end{tabular}
\end{table}

While the workaround eliminates the step-function artifact, it does
not reproduce the original's \texttt{tsd} values.
Table~\ref{tab:boundary_three} shows that the workaround
($t_{\text{sd}} = 0.5$) produces systematically higher boundary
carryover than the original ($t_{\text{sd}} = 1$) at every half-life.
The correct fix is to adopt the original's interval-based accumulation
logic (Section~\ref{sec:fix}).

\section{Open Question: The Original's Power Behavior}
\label{sec:open}

Independent of the 2025 \texttt{tsd} bug, the published Hendrickson
heatmaps \cite{hendrickson2020} show power declining with increasing
carryover half-life. Whether this decline is steeper than
pharmacologically expected is the motivating question for the
\texttt{carryover\_focus/} analysis. Once the \texttt{tsd} bug is
fixed, any remaining excessive power dropoff must have a different
mechanism. Three candidates merit investigation.

\subsection{Scale Factor Mismatch Between DGP and Analysis}

The original code uses different scale factors for carryover in the
data-generating process and the analysis model:

\begin{itemize}[nosep]
  \item \textbf{DGP} (\texttt{generateData.R}, line 96):
    \texttt{scalefactor = 2} (default).
    \[
      \text{carryover}_{\text{DGP}} =
        \text{br}_{\text{prev}} \times
        \left(\tfrac{1}{2}\right)^{2 \cdot t_{\text{sd}} / t_{1/2}}
    \]

  \item \textbf{Analysis} (\texttt{lme\_analysis.R}, line 109):
    \texttt{carryover\_scalefactor = 1} (default).
    \[
      \texttt{Dbc} =
        \left(\tfrac{1}{2}\right)^{1 \cdot t_{\text{sd}} / t_{1/2}}
    \]
\end{itemize}

The DGP generates faster decay (higher effective exponent) than the
analysis model assumes. At $t_{\text{sd}} = 1$, $t_{1/2} = 0.2$:

\begin{itemize}[nosep]
  \item DGP carryover: $(1/2)^{2 \times 1/0.2} = (1/2)^{10}
    \approx 0.001$
  \item Analysis \texttt{Dbc}: $(1/2)^{1 \times 1/0.2} = (1/2)^{5}
    = 0.031$
\end{itemize}

The analysis model overestimates the residual drug effect by a factor
of $\sim$30 at this timepoint. This systematic misspecification could
contribute to power loss by distorting the interaction contrast that
Model~A attempts to estimate.

\subsection{Cumulative Carryover in \texttt{generateData()}}

In the DGP (\texttt{generateData.R}, lines 92--99), carryover is
added to \texttt{brmeans} cumulatively:

\begin{verbatim}
for (p in 2:nP) {
  if (!d[p]$onDrug) {
    if (d[p]$tsd > 0) {
      brmeans[p] <- brmeans[p] +
        brmeans[p-1] * (1/2)^(sf * d$tsd[p] / t_half)
    }
  }
}
\end{verbatim}

This recursive structure means the biological response at off-drug
timepoints depends on the entire preceding on-drug trajectory, not
just the value at discontinuation. The interaction with the Gompertz
response curve (\texttt{modgompertz()}) and the multi-factor
correlation structure may produce nonlinear contamination patterns
that the analysis model cannot fully capture.

\subsection{Correlation Structure Interaction}

The original DGP generates correlated random components across
timepoints via a multivariate normal draw with a structured
correlation matrix (autocorrelation within factors, cross-correlation
between factors). The carryover adjustment modifies the mean structure
but not the correlation matrix. As $t_{1/2}$ increases, the mismatch
between the adjusted means and the original correlation structure
grows, potentially producing model misspecification that reduces the
interaction estimate.

\subsection{Status}

These three candidates are not mutually exclusive. Determining their
relative contributions requires running the corrected 2025 simulation
(with original-matching \texttt{tsd}) and comparing its power curves
against the published results. The scale factor mismatch
(Section~6.1) is the most straightforward to test: it can be isolated
by running simulations with matched scale factors (both set to~1 or
both set to~2) and observing whether the power curve changes shape.

\section{Recommended Path Forward}
\label{sec:fix}

\begin{enumerate}
  \item \textbf{Fix \texttt{tsd} in
    \texttt{carryover\_power\_simulation.R}.} Replace the
    calendar-subtraction logic with the original's interval-based
    accumulation. Verify that the corrected code produces identical
    \texttt{tsd} values to \texttt{buildtrialdesign()} for all four
    paths (Table~\ref{tab:orig_tsd}).

  \item \textbf{Re-run the two-scenario comparison} with the corrected
    \texttt{tsd} and \texttt{tsd\_min}~$= 0$ to establish a faithful
    reproduction of Hendrickson's simulation as a baseline.

  \item \textbf{Compare power curves} from the corrected simulation
    against the published heatmaps. If power still drops more rapidly
    than expected, investigate the scale factor mismatch by running
    with matched scale factors.

  \item \textbf{Evaluate \texttt{tsd\_min} as a design modification.}
    If the original's power behavior proves excessive, a minimum
    washout interval represents a pharmacologically defensible
    improvement: in clinical practice, there is always a nonzero gap
    between the last dose and the next assessment. This should be
    evaluated as a deliberate modeling choice with its own simulation
    results, not as a bug fix.
\end{enumerate}

\section{Conclusions}

\begin{enumerate}
  \item \textbf{The step-function power cliff in the 2025 simulation
    is caused by a \texttt{tsd} computation bug.} The 2025 code
    computes $t_{\text{sd}} = \text{week} - \text{discontinuation\_week}$,
    producing $t_{\text{sd}} = 0$ at boundary measurements. The
    original computes \texttt{tsd} as cumulative off-drug interval
    widths, producing $t_{\text{sd}} \geq 1$ at all off-drug
    timepoints. The step function does not exist in the original
    framework.

  \item \textbf{The \texttt{tsd\_min} workaround partially
    compensates but does not match the original.} With
    $t_{\text{sd,min}} = 0.5$, the boundary carryover at
    $t_{1/2} = 0.1$ is 0.031; with the original \texttt{tsd}, it is
    0.001. The correct fix is to adopt the original's interval-based
    accumulation.

  \item \textbf{The original's power behavior is a separate, open
    question.} The published heatmaps show power declining with
    increasing $t_{1/2}$. Whether this decline is excessive requires
    investigation after the \texttt{tsd} bug is fixed. Leading
    candidates include the scale factor mismatch between the DGP
    (\texttt{scalefactor}~$= 2$) and the analysis model
    (\texttt{carryover\_scalefactor}~$= 1$), and the cumulative
    carryover structure in \texttt{generateData()}.

  \item \textbf{The two-scenario comparison framework remains valid.}
    Comparing Hendrickson's \texttt{Dbc} predictor against the naive
    \texttt{treatment} indicator is the right experimental design.
    The simulation results in
    Tables~\ref{tab:uncorrected}--\ref{tab:corrected} document the
    behavior of the buggy code and its workaround; new results with
    the corrected \texttt{tsd} are needed.
\end{enumerate}

\section*{References}

\begingroup
\renewcommand{\section}[2]{}
\begin{thebibliography}{9}

\bibitem{hendrickson2020}
Hendrickson RC, Thomas RG, Schork NJ.
Optimizing aggregated N-of-1 trial designs for predictive biomarker
validation: Statistical methods and theoretical findings.
\textit{Frontiers in Digital Health}. 2020;2:13.
doi:~\href{https://doi.org/10.3389/fdgth.2020.00013}%
  {10.3389/fdgth.2020.00013}

\bibitem{wang2019}
Wang Y, Schork NJ.
Power and design issues in crossover-based N-of-1 clinical trials with
fixed data collection periods.
\textit{Healthcare}. 2019;7(3):84.
doi:~\href{https://doi.org/10.3390/healthcare7030084}%
  {10.3390/healthcare7030084}

\bibitem{yap2024}
Yap C, Wang X, Harhay MO.
Behavioral carry-over effect and power consideration in crossover
trials.
\textit{Biometrics}. 2024;80(2):ujae023.
doi:~\href{https://doi.org/10.1093/biomtcs/ujae023}%
  {10.1093/biomtcs/ujae023}

\bibitem{blackston2019}
Blackston JW, Chapple AG, McGree JM, McDonald S, Nikles J.
Comparison of aggregated N-of-1 trials with parallel and crossover
randomized controlled trials using simulation studies.
\textit{Healthcare}. 2019;7(4):137.
doi:~\href{https://doi.org/10.3390/healthcare7040137}%
  {10.3390/healthcare7040137}

\bibitem{senn2019}
Senn S.
Sample size considerations for n-of-1 trials.
\textit{Statistical Methods in Medical Research}. 2019;28(2):372--383.
doi:~\href{https://doi.org/10.1177/0962280217726801}%
  {10.1177/0962280217726801}

\end{thebibliography}
\endgroup

\bigskip
\hrule
\medskip
{\small\itshape
  Rendered on 2026-03-18 at 09:37 PDT.\\
  Source: \texttt{\~/prj/alz/01-pmsimstats/pmsimstats2025/%
    analysis/scripts/carryover\_focus/predictor\_comparison.tex}
}

\end{document}
